% !TeX root = ../../Book1.tex
\section{绝对连续函数与微积分基本定理}
\label{b1:sec:1604}

绝对连续性把“小区间族的总长度很小”转化为“函数增量的总和很小”。它比一致连续性强，又允许导数在局部无界；更重要的是，它恰好给出可以由可积导数恢复的实函数。本节的结果是后续一维 Sobolev 代表理论的直接前置，而不属于可以跳过的度量曲线选读内容。

% 实读来源：Fremlin2016Measure，作者官方 chap22.pdf，225B--E（该文件印刷 pp.27--29）。
% 积分表示以变差、单调分解、广义逆推前和 RN 补写，不把未写的 Stieltjes 构造当作先修。

\subsection{绝对连续性与全变差}
\label{b1:subsec:160401}

\begin{definition}
\label{b1:def:absolutely-continuous-function}
设 \(f:[a,b]\to\mathbb R\)，其中 \(a<b\)。若对每个 \(\varepsilon>0\)，存在 \(\delta>0\)，使任意有限族内点互不相交的区间 \([s_i,t_i]\subseteq[a,b]\) 满足
\[
\sum_i(t_i-s_i)<\delta
\quad\Longrightarrow\quad
\sum_i\bigl|f(t_i)-f(s_i)\bigr|<\varepsilon,
\]
则称 \(f\) 为\term[juedui-lianxu-hanshu]{绝对连续函数}[absolutely continuous function]。
\end{definition}

这里只取一个区间，便可看出绝对连续函数一致连续。Lipschitz 函数满足定义，因为增量绝对值之和不超过 Lipschitz 常数乘以区间长度之和；反向一般不成立，例如本节末的平方根函数。

\begin{lemma}
\label{b1:lem:ac-real-variation}
设 \(f:[a,b]\to\mathbb R\) 绝对连续，定义
\[
V(t)=\sup\Bigl\{\sum_{i=1}^n
  \bigl|f(t_i)-f(t_{i-1})\bigr|:
a=t_0<\cdots<t_n=t\Bigr\},
\quad V(a)=0.
\]
则 \(V\) 有限、单调不减且绝对连续。对 \(a\leq s<t\leq b\)，有
\[
\bigl|f(t)-f(s)\bigr|\leq V(t)-V(s),
\]
而右端正是 \(f\) 在子区间 \([s,t]\) 上的全变差。
\end{lemma}
\begin{proof}
在一个分割中插入中间点，由三角不等式，增量绝对值之和不会减小；反过来，可以拼接两个子区间的分割。这证明分割上确界对子区间可加，允许暂取无穷值。

对 \(\varepsilon=1\) 取绝对连续性定义中的 \(\delta\)，把 \([a,b]\) 分成有限个长度小于 \(\delta\) 的区间。在其中任一区间内，任意分割的各小段总长度小于 \(\delta\)，所以增量绝对值之和小于 \(1\)。把一个任意分割加上这些预定端点，便得到总上界，故 \(V(b)<\infty\)。单调性、子区间公式及增量估计随之成立。

最后，给定 \(\varepsilon>0\)，对绝对连续性取使增量绝对值之和小于 \(\varepsilon/2\) 的 \(\delta\)。若一族不交区间的总长度小于 \(\delta\)，在每个区间内部选分割来逼近其分割上确界。所有分割小段仍组成总长度小于 \(\delta\) 的不交族。因此，取这些有限个上确界后，
\[
\sum_i\bigl(V(t_i)-V(s_i)\bigr)\leq\varepsilon/2<\varepsilon.
\]
这正是 \(V\) 的绝对连续性。
\end{proof}

\subsection{微积分基本定理}
\label{b1:subsec:160402}

Fremlin 的《Measure Theory》\cite[225B--E]{Fremlin2016Measure}系统讨论绝对连续函数及其积分表示。下面采用\chapref{b1:ch:06}的 Radon--Nikodym 定理来证明表示方向，使这里不依赖尚未建立的一般有界变差函数微分理论。

\begin{theorem}[微积分基本定理，绝对连续函数]
\label{b1:thm:ac-fundamental-calculus}
实函数 \(f:[a,b]\to\mathbb R\) 绝对连续，当且仅当存在 \(g\in L^1([a,b])\)，使
\[
f(t)=f(a)+\int_a^t g(r)\,\mathrm dr
\quad(t\in[a,b]).
\]
在这种情形下，\(f'=g\) 几乎处处，因而可以取 \(g=f'\)。
\end{theorem}
证明的正向先以全变差函数把一般 \(f\) 分解为两个单调绝对连续函数：
\[
 f\longrightarrow V\longrightarrow P,Q.
\]
对其中任一单调函数 \(F\)，再用广义逆 \(q\) 把区间上的 Lebesgue 测度推前为增量测度
\(\mu_F\)，由绝对连续性证明 \(\mu_F\ll m_1\)，最后应用 Radon--Nikodym 定理取得积分密度。
反向则直接用积分的绝对连续性。因而正向证明的主链是
\[
 F\longrightarrow q\longrightarrow\mu_F
 \longrightarrow\mu_F\ll m_1
 \longrightarrow\frac{\mathrm d\mu_F}{\mathrm dm_1}.
\]
\begin{proof}
先设 \(f\) 绝对连续，令 \(V\) 为\cref{b1:lem:ac-real-variation}中的全变差函数。两个函数
\[
P(t)=\frac{V(t)+f(t)-f(a)}2,
\quad
Q(t)=\frac{V(t)-f(t)+f(a)}2
\]
都单调不减，因为 \(V(t)-V(s)\geq|f(t)-f(s)|\)。它们还绝对连续，在 \(a\) 处为零，且 \(f-f(a)=P-Q\)。所以只需为任意单调不减的绝对连续函数 \(F\)，在 \(F(a)=0\) 的归一化下，构造积分密度。

令 \(M=F(b)\)。若 \(M=0\)，则 \(F=0\)。若 \(M>0\)，对 \(0<u<M\) 定义
\[
q(u)=\min\{t\in[a,b]:F(t)\geq u\}.
\]
这个最小值由 \(F\) 的连续性与单调性存在，并且
\[
q(u)\leq t\quad\Longleftrightarrow\quad u\leq F(t).
\]
所以 \(q\) 是 Borel 可测函数。把 \((0,M)\) 上的 Lebesgue 测度经 \(q\) 推前，得到 \([a,b]\) 上的有限正 Borel 测度 \(\mu_F\)，满足
\[
\mu_F([a,t])=F(t),
\quad
\mu_F((s,t])=F(t)-F(s).
\]
又由 \(F\) 连续，这个测度没有原子，包括两个端点。

现证明 \(\mu_F\ll m_1\)。设 Borel 集 \(N\subseteq[a,b]\) 的 Lebesgue 测度为零，给定 \(\varepsilon>0\)，取 \(F\) 的绝对连续性所对应的 \(\delta\)。由 Lebesgue 测度的外正则性，存在开集 \(O\supseteq N\)，使 \(m_1(O)<\delta\)。开集 \(O\) 与 \([a,b]\) 的交可分成可数个内点互不相交的区间；端点有没有包含不影响 \(\mu_F\) 的值。每个有限子族的 \(\mu_F\) 测度之和，就是相应的 \(F\) 增量之和，不超过 \(\varepsilon\)。由可数可加性，\(\mu_F(N)\leq\mu_F(O\cap[a,b])\leq\varepsilon\)。令 \(\varepsilon\downarrow0\)，得到绝对连续性。

由\cref{b1:thm:radon-nikodym-finite}，存在非负 \(g_F\in L^1([a,b])\)，使 \(\mu_F=g_Fm_1\)，因此 \(F(t)=\int_a^t g_F\)。分别用于 \(P,Q\)，便得到所需 \(g=g_P-g_Q\)。这同时补出了本证明所需的单调函数生成测度的步骤。

反过来，若 \(f(t)=f(a)+\int_a^t g\)，则对任意不交区间族，
\[
\sum_i|f(t_i)-f(s_i)|
\leq\int_{\bigcup_i[s_i,t_i]}|g(r)|\,\mathrm dr.
\]
积分的绝对连续性，即\cref{b1:prop:integral-absolute-continuity}，证明 \(f\) 绝对连续。最后，在 \(g\) 的每个内部 Lebesgue 点 \(t\)，
\[
\left|\frac{f(t+h)-f(t)}h-g(t)\right|
\leq\frac1{|h|}\int_{t-|h|}^{t+|h|}|g(r)-g(t)|\,\mathrm dr
\longrightarrow0.
\]
\cref{b1:thm:lebesgue-point}遂给出 \(f'=g\) 几乎处处。
\end{proof}

例如，\(t\mapsto\sqrt t\) 在 \([0,1]\) 上绝对连续，因为它是 \((2\sqrt t)^{-1}\) 的积分，但它在零点附近不是 Lipschitz 的。另一方向的边界由\cref{b1:ex:14-cantor-function}中的 Cantor 函数给出：它连续、单调且导数几乎处处为零，却不是常数，所以不是绝对连续函数。仅有连续性和几乎处处可微性，还不足以保证微积分基本定理。

有限区间上的正则性条件决定导数能否恢复全部函数增量，见\cref{b1:tab:visual-t09}。
% Source fingerprint: 89a41d06e175d9aa5ca92aa3ddd8667a41bbec1369852f31d74b225a69b2e5b8
% T09: raw TeX cells preserved; no numeric highlighting.
\begin{table}[htbp]
\centering\small\renewcommand{\arraystretch}{1.18}\setlength{\tabcolsep}{4pt}
\caption[Lipschitz、绝对连续与有界变差]{Lipschitz、绝对连续与有界变差。这里的一维变差按划分定义。第27章从分布导数测度重新组织该理论；函数代表与端点取值始终按正文约定。}
\label{b1:tab:visual-t09}
\begin{tabularx}{\linewidth}{@{}>{\raggedright\arraybackslash}p{3.1cm}>{\raggedright\arraybackslash}p{4.7cm}>{\raggedright\arraybackslash}X@{}}
\toprule
有限区间上的函数类 & 导数与积分恢复 & 不能逆推的例子 \\
\midrule
Lipschitz & 绝对连续；导数几乎处处存在且本性有界；积分恢复函数增量 & 绝对连续函数不必 Lipschitz，如平方根 \\
绝对连续 & 导数属于 \(L^1\)，且 \(f(t)-f(a)=\int_a^t f^{\prime}(s)\,\mathrm{d}s\) & 仅有连续性和几乎处处可微不足以保证积分恢复 \\
有界变差 & 划分增量绝对值之和有统一上界；导数可积，但可有奇异变化 & Cantor 函数连续、几乎处处导数为零，却不是常值 \\
一般连续函数 & 不保证有界变差或几乎处处可微 & 本章三角峰例子可使划分增量和发散 \\
\bottomrule
\end{tabularx}
\end{table}

